% paper3_dark_matter.tex
% Observer-Dependent Dark Matter: Gravitational Entanglement
% and the MOND Scale from Quantum Retrodiction
% For submission to Physical Review D
% Author: Sheng-Kai Huang
% Compile with: pdflatex paper3_dark_matter.tex

\documentclass[prd,twocolumn,superscriptaddress,longbibliography,floatfix]{revtex4-2}

\usepackage{amsmath}
\usepackage{amssymb}
\usepackage{amsfonts}
\usepackage{amsthm}
\usepackage{bm}
\usepackage{hyperref}
\usepackage{booktabs}
\usepackage{graphicx}

\newtheorem{theorem}{Theorem}
\newtheorem{corollary}[theorem]{Corollary}
\newtheorem{conjecture}[theorem]{Conjecture}
\newtheorem*{proposition}{Proposition}
\newtheorem{definition}{Definition}

% Convenience macros (consistent with Papers 1, 2, 5)
\newcommand{\kB}{k_{\mathrm{B}}}
\newcommand{\Tr}{\mathrm{Tr}}
\newcommand{\id}{\mathrm{id}}
\newcommand{\Petz}{\mathcal{R}}
\newcommand{\chan}{\mathcal{N}}
\newcommand{\hilb}{\mathcal{H}}
\newcommand{\code}{\mathcal{C}}
\newcommand{\KL}{D}
\newcommand{\ket}[1]{|#1\rangle}
\newcommand{\bra}[1]{\langle#1|}
\newcommand{\braket}[2]{\langle#1|#2\rangle}
\newcommand{\op}[2]{|#1\rangle\langle#2|}
\newcommand{\abs}[1]{\left|#1\right|}
\newcommand{\norm}[1]{\left\|#1\right\|}
\newcommand{\tauasym}{\tau}
\newcommand{\algO}{\mathcal{A}}
\newcommand{\Msun}{M_\odot}

\begin{document}

\title{Observer-Dependent Dark Matter: Gravitational Entanglement\\
and the MOND Scale from Quantum Retrodiction}

\author{Sheng-Kai Huang}
\email{akai@fawstudio.com}
\affiliation{Independent Researcher}

\date{\today}

\begin{abstract}
We propose that dark matter phenomenology at galactic scales
arises from the observer-dependence of gravitational entropy
production $\Sigma$.
A classical observer, with access only to electromagnetic
degrees of freedom, assigns a larger entropy production
$\Sigma_{\mathrm{cl}}$ to a gravitational process than a
quantum observer who can additionally access the gravitational
entanglement structure: $\Sigma_{\mathrm{cl}} \geq
\Sigma_{\mathrm{q}}$, by the data processing inequality.
The difference $\Sigma_{\mathrm{DM}} \equiv
\Sigma_{\mathrm{cl}} - \Sigma_{\mathrm{q}}$ manifests as
apparent dark matter.
We develop this framework quantitatively by connecting
observer-dependent $\tauasym = 1 - F$ (Paper~V) to the
scale-dependent gravitational channel: the anomalous dimension
$\eta = 1$ identified by Kumar [arXiv:2509.05246], reproduced
within the $\tauasym$ framework via the data processing
inequality and extensivity, produces a logarithmic correction
to the Newtonian potential, yielding flat rotation curves with
a universal crossover scale $k_* \approx 2.7 \times 10^{-2}
\;\mathrm{kpc}^{-1}$, consistent with the 100-galaxy SPARC
validation of Gubitosi \emph{et al.}
The MOND acceleration scale $a_0 \approx
cH_0/(2\pi)$ is related to the KMS thermal periodicity of the
de~Sitter vacuum through the Crooks fluctuation theorem,
interpreting the MOND transition as the onset of marginal
irreversibility ($\Sigma \sim 1$).
A Crooks-inspired interpolation function
$\mu(x) = 1 - e^{-x}$ with $x = a/a_0$ is proposed and
shown to be competitive with standard MOND interpolations.
The Bullet Cluster is addressed through an entanglement
rigidity argument.
The Khronon coupling
$\mu_0 = H_0/c$---the unique dimensionally-determined value
from the de~Sitter temperature---with anomalous dimension
$\eta_\mu = 1$ produces dust-like behavior
($\tilde{w}_0 \sim 4 \times 10^{-10}$) at CMB scales,
providing a candidate for the pressureless mode needed for
the acoustic peaks.
The exact $\Omega_{\mathrm{DM}} h^2$ prefactor requires
the numerical computation of Paper~IV, and this CMB
resolution remains unverified.
\end{abstract}

\maketitle

%=======================================================================
\section{Introduction}
\label{sec:intro}
%=======================================================================

The dark matter problem is one of the central puzzles in
physics.  Galaxy rotation curves flatten at large
radii~\cite{Rubin1980}, the cosmic microwave background (CMB)
acoustic peaks require a pressureless, non-baryonic component
with $\Omega_{\mathrm{CDM}} h^2 \approx 0.12$~\cite{Planck2018},
and gravitational lensing in the Bullet Cluster shows mass
displaced from the dominant baryonic component~\cite{Clowe2006}.
These observations are well described by the
$\Lambda$CDM concordance model, in which cold dark matter
(CDM) particles constitute roughly 27\% of the energy budget
of the universe.

Yet four decades of direct detection experiments have found no
CDM particles~\cite{Schumann2019}.  This null result, combined
with the tight empirical scaling relations that CDM does not
predict---the radial acceleration relation
(RAR)~\cite{McGaughRAR2016}, the baryonic Tully--Fisher
relation~\cite{McGaughBTFR2000}, and the near-universal
acceleration scale $a_0 \approx 1.2 \times 10^{-10}\;
\mathrm{m/s}^2$---motivates the search for alternative
explanations.

Modified Newtonian dynamics (MOND)~\cite{Milgrom1983}
reproduces galactic phenomenology with a single parameter
$a_0$ but lacks a satisfactory relativistic completion
and fails at CMB scales~\cite{Skordis2021}.
Verlinde's emergent gravity~\cite{Verlinde2016} derives
$a_0 \sim cH_0/6$ from volume-law entanglement in
de~Sitter space and achieves zero-parameter fits to
rotation curves, but has no covariant formulation and
fails quantitatively at cluster scales~\cite{Ettori2019}.
The AeST theory of Skordis and
Zlo\v{s}nik~\cite{Skordis2021} fits both galactic rotation
curves and the CMB power spectrum, but does so by introducing
a vector field and scalar field that are, in a precise sense,
functionally equivalent to a dark sector.

In this paper, we propose a different perspective.  We argue
that ``dark matter'' at galactic scales is the gravitational
information that classical observers cannot access.  The
framework rests on the observer-dependent temporal asymmetry
parameter $\tauasym_O = 1 - F(\rho, \Petz \circ \chan_O(\rho))$
developed in Ref.~\cite{Paper5} (hereafter Paper~V), where $F$
is the fidelity of Petz recovery and $\chan_O$ is the effective
channel determined by the observer's accessible
algebra~$\algO_O$.  A classical astronomer, restricted to
electromagnetic observables ($\algO_{\mathrm{cl}}$), assigns a
larger entropy production $\Sigma_{\mathrm{cl}}$ to the same
gravitational process than a hypothetical quantum observer with
access to the full gravitational entanglement structure
($\algO_{\mathrm{q}} \supseteq \algO_{\mathrm{cl}}$).  The
data processing inequality (DPI) guarantees
$\Sigma_{\mathrm{cl}} \geq \Sigma_{\mathrm{q}}$, and the
excess $\Sigma_{\mathrm{DM}} = \Sigma_{\mathrm{cl}} -
\Sigma_{\mathrm{q}}$ manifests as apparent dark matter.

The paper is organized as follows.
Section~\ref{sec:framework} formalizes observer-dependent dark
matter using the subalgebra framework of Paper~V.
Section~\ref{sec:galactic} develops the galactic-scale
phenomenology: running Newton's constant from the
renormalization group (RG) flow of the gravitational channel,
flat rotation curves, the RAR, and the baryonic Tully--Fisher
relation.
Section~\ref{sec:a0} relates the MOND acceleration scale
to the Crooks fluctuation theorem in de~Sitter space and
proposes a Crooks-inspired interpolation function.
Section~\ref{sec:bullet} addresses the Bullet Cluster.
Section~\ref{sec:cosmo} discusses cosmic acceleration and
dark energy.
Section~\ref{sec:cmb} assesses the CMB gap and
identifies the structural parallel with AeST/Khronon theories;
the running Khronon coupling $\mu_0 = H_0/c$ with
$\eta_\mu = 1$ then provides a resolution yielding dust-like
behavior at CMB scales, with the full treatment deferred to
Paper~IV.
Section~\ref{sec:predictions} lists testable predictions,
and Sec.~\ref{sec:discussion} provides a critical assessment
of what is genuinely new and what remains open.

%=======================================================================
\section{Framework: Observer-Dependent Dark Matter}
\label{sec:framework}
%=======================================================================

\subsection{Review: observer-dependent $\tauasym$}

In Paper~V~\cite{Paper5}, each observer $O$ is identified with
a von~Neumann subalgebra $\algO_O \subseteq
\algO_{\mathrm{total}}$ of the total algebra of observables.
The observer's effective channel is the conditional expectation
$\chan_O = \Tr_{\algO_O^c}$, and the observer-dependent
temporal asymmetry is
\begin{equation}
\tauasym_O = 1 - F\!\big(\rho,\,
\tilde{\Petz}_{\sigma,\chan_O}(\chan_O(\rho))\big),
\label{eq:tau_O}
\end{equation}
bounded by the JRSWW inequality~\cite{Junge2018}:
\begin{equation}
\tauasym_O \leq 1 - \exp(-\Sigma_O/2),
\label{eq:tau_bound}
\end{equation}
where $\Sigma_O = \KL(\rho\|\sigma) -
\KL(\chan_O(\rho)\|\chan_O(\sigma))$ is the relative entropy
lost by restricting to $\algO_O$.  The monotonicity theorem
(Theorem~1 of Paper~V) states: if $O_1 \leq O_2$ (i.e.,
$\algO_{O_1} \subseteq \algO_{O_2}$), then
$\tauasym_{O_1} \geq \tauasym_{O_2}$ and
$\Sigma_{O_1} \geq \Sigma_{O_2}$.

\subsection{Classical and quantum observers for a galaxy}

\begin{definition}[Classical observer]
\label{def:classical}
The classical galactic observer $O_{\mathrm{cl}}$ has
accessible algebra
\begin{equation}
\algO_{\mathrm{cl}} = \algO_{\mathrm{EM}} \otimes
\mathbf{1}_{\mathrm{grav}} \otimes \mathbf{1}_{\mathrm{env}},
\end{equation}
where $\algO_{\mathrm{EM}}$ contains all observables
constructible from electromagnetic field modes and their
couplings to baryonic matter (photometry, spectroscopy,
21\,cm emission, thermal bremsstrahlung).
\end{definition}

\begin{definition}[Quantum observer]
\label{def:quantum}
The quantum galactic observer $O_{\mathrm{q}}$ has
\begin{equation}
\algO_{\mathrm{q}} = \algO_{\mathrm{EM}} \vee
\algO_{\mathrm{grav}},
\end{equation}
the algebra generated by $\algO_{\mathrm{EM}}$ and the
gravitational entanglement observables
$\algO_{\mathrm{grav}}$, which include vacuum entanglement
between spatial regions of the gravitational field,
volume-law entanglement from the de~Sitter
background~\cite{Verlinde2016}, and modular flow observables
of the gravitational field~\cite{DorauMuch2025}.
\end{definition}

Since $\algO_{\mathrm{cl}} \subset \algO_{\mathrm{q}}$, the
monotonicity theorem gives:

\begin{theorem}[Dark matter inequality]
\label{thm:dm}
For any gravitational process with reference state
$\sigma$,
\begin{equation}
\Sigma_{\mathrm{cl}} \geq \Sigma_{\mathrm{q}} \geq 0,
\label{eq:dm_ineq}
\end{equation}
with equality if and only if the gravitational entanglement
carries no information beyond what is accessible
electromagnetically.
\end{theorem}

\begin{proof}
Immediate from the DPI applied to the conditional expectation
$\mathcal{E}: \algO_{\mathrm{q}} \to \algO_{\mathrm{cl}}$,
since $\chan_{\mathrm{cl}} = \mathcal{E} \circ
\chan_{\mathrm{q}}$.
\end{proof}

\begin{definition}[Dark matter entropy production]
\label{def:sigma_dm}
The observer-dependent dark matter entropy production is
\begin{equation}
\Sigma_{\mathrm{DM}} \equiv \Sigma_{\mathrm{cl}} -
\Sigma_{\mathrm{q}} \geq 0.
\label{eq:sigma_dm}
\end{equation}
\end{definition}

\subsection{From $\Sigma_{\mathrm{DM}}$ to apparent dark
matter mass}

Using the connection $\Sigma_{\mathrm{grav}} = -\ln(-g_{00})$
established in Ref.~\cite{Paper2} (Paper~II), the gravitational
potential is
$\Phi(r) = -(c^2/2)\,\Sigma(r)$ in the weak-field limit.
The mass inferred by the classical observer within radius $r$
is
\begin{equation}
M_{\mathrm{obs}}(r) = \frac{c^2 r^2}{2G_N}\,
\Sigma_{\mathrm{cl}}'(r),
\label{eq:M_obs}
\end{equation}
and the baryonic mass attributed by the quantum observer is
\begin{equation}
M_{\mathrm{bar}}(r) = \frac{c^2 r^2}{2G_N}\,
\Sigma_{\mathrm{q}}'(r).
\label{eq:M_bar}
\end{equation}
The apparent dark matter mass is therefore
\begin{equation}
M_{\mathrm{DM}}(r) = M_{\mathrm{obs}}(r) - M_{\mathrm{bar}}(r)
= \frac{c^2 r^2}{2G_N}\,\Sigma_{\mathrm{DM}}'(r).
\label{eq:M_DM}
\end{equation}

The challenge is to determine $\Sigma_{\mathrm{DM}}(r)$.
We now show that the scale-dependent gravitational channel
provides the answer.

%=======================================================================
\section{Galactic Scale: Rotation Curves}
\label{sec:galactic}
%=======================================================================

\subsection{Scale-dependent gravitational channel}

The gravitational field acts as a quantum channel
$\chan_{\mathrm{grav}}$ that processes quantum information.
At different momentum scales $k$ (corresponding to distance
scales $r \sim 1/k$), the channel has different properties
because Newton's constant $G$ runs under the renormalization
group~\cite{Reuter2004,Donoghue1994}.  The entropy production
at scale $k$ is
\begin{equation}
\Sigma(k) = \KL(\rho\|\sigma) -
\KL\!\big(\chan_{\mathrm{grav}}(k)(\rho)
\big\|\chan_{\mathrm{grav}}(k)(\sigma)\big).
\label{eq:sigma_k}
\end{equation}

The DPI, applied to the RG flow---which is itself a quantum
channel~\cite{CasiniHuerta2017}---gives a monotonicity
condition:
\begin{equation}
\Sigma(k_{\mathrm{IR}}) \geq \Sigma(k_{\mathrm{UV}}) \geq 0.
\label{eq:sigma_mono}
\end{equation}
The gravitational channel at infrared (IR) scales is
noisier---it loses more information---than at ultraviolet (UV)
scales.  This monotonicity is the entropic $c$-theorem
of Casini and Huerta~\cite{CasiniHuerta2012}, reinterpreted
as a statement about the gravitational channel.

\subsection{Running $G$ and the logarithmic correction}

The anomalous dimension $\eta$ governs the IR running:
\begin{equation}
\frac{d\ln G}{d\ln k} = -\eta.
\label{eq:anomalous}
\end{equation}
For $k \ll k_*$, the solution is
$G(k) = G_N(k_*/k)^\eta$.
Kumar~\cite{Kumar2025} showed from QFT first principles that
$\eta = 1$ is the unique marginal value: it is the boundary
between corrections that die off ($\eta < 1$) and corrections
that grow without bound ($\eta > 1$).  Within the $\tauasym$
framework, this value is independently reproduced by the DPI
plus extensivity of de~Sitter entanglement entropy
(see Sec.~\ref{sec:discussion} for a sketch of the argument;
the complete proof requires the perturbative RG analysis
of Paper~IV~\cite{Paper4}).  In position space, $\eta = 1$ produces the
unique scale-invariant, rotationally symmetric deformation
of the Newtonian potential:
\begin{equation}
G(r) = G_N\!\left[1 + \frac{2k_* r}{\pi}\right],
\quad r \gg 1/k_*,
\label{eq:G_running}
\end{equation}
with the coefficient $2/\pi$ being universal and
regulator-independent.

The modified potential, circular velocity, and effective
entropy production are:
\begin{align}
\Phi(r) &= -\frac{G_N M}{r} +
\frac{2G_N M k_*}{\pi}\ln\!\left(\frac{r}{r_0}\right),
\label{eq:phi}\\
v^2(r) &= \frac{G_N M}{r} +
\frac{2G_N M k_*}{\pi},
\label{eq:v2}\\
\Sigma_{\mathrm{grav}}(r) &= \frac{r_s}{r} +
\frac{2k_* r_s}{\pi}\ln\!\left(\frac{r}{r_0}\right),
\label{eq:sigma_grav}
\end{align}
where $r_s = 2G_N M/c^2$ is the Schwarzschild radius and
$r_0$ is a reference scale.  Equation~\eqref{eq:v2} yields
a flat rotation curve at $r \gg 1/k_*$:
\begin{equation}
V_0^2 = \frac{2G_N M_{\mathrm{bar}}\, k_*}{\pi},
\label{eq:V0}
\end{equation}
which is constant.

The logical chain connecting established results is:
\begin{align}
&\text{RG flow is a quantum channel}
\;\cite{CasiniHuerta2017}
\nonumber\\
\xrightarrow{\mathrm{DPI}}\;&
\text{QRE monotonicity under RG}
\nonumber\\
\xrightarrow{\mathrm{Dorau\text{-}Much}}\;&
\text{QRE} \to \text{semiclassical Einstein eqns.}
\;\cite{DorauMuch2025}
\nonumber\\
\xrightarrow{\mathrm{Kumar}}\;&
\text{Running }G\text{ with }\eta=1
\to \text{flat rotation curves}
\;\cite{Kumar2025}.
\nonumber
\end{align}

\subsection{SPARC galaxies and the universal crossover scale}

Gubitosi, Piattella, and Casarini~\cite{Gubitosi2024}
validated the running-$G$ framework (termed RGGR) against 100
galaxies from the SPARC database~\cite{SPARC2016}, finding it
competitive with NFW dark matter halos (RGGR preferred for
47 galaxies, NFW for 40, inconclusive for 13 by the BIC
criterion).

Kumar~\cite{Kumar2025} fitted three SPARC galaxies and found a
consistent crossover scale
(Table~\ref{tab:sparc}):

\begin{table}[t]
\caption{\label{tab:sparc}Crossover scale $k_*$ for three
SPARC galaxies from Ref.~\cite{Kumar2025}.  The near-universality
of $k_*$ across a factor of $\sim\!3$ in baryonic mass is a
specific prediction of the framework.}
\begin{ruledtabular}
\begin{tabular}{lccc}
Galaxy & $M_{\mathrm{bar}}\;[\Msun]$ & $V_0\;[\mathrm{km/s}]$
& $k_*\;[\mathrm{kpc}^{-1}]$ \\
\colrule
S:700624 & $1.01\times 10^{12}$ & $271 \pm 3$
& $(2.66 \pm 0.05)\!\times\!10^{-2}$ \\
S:700916 & $5.85\times 10^{11}$ & $210 \pm 2$
& $(2.76 \pm 0.06)\!\times\!10^{-2}$ \\
S:705253 & $3.55\times 10^{11}$ & $159 \pm 2$
& $(2.60 \pm 0.06)\!\times\!10^{-2}$ \\
\end{tabular}
\end{ruledtabular}
\end{table}

The universal $k_* \approx 2.7 \times 10^{-2}\;\mathrm{kpc}^{-1}$
(crossover radius $r_c = 1/k_* \approx 37\;\mathrm{kpc}$) is
a specific prediction absent from NFW halo models, which
require two parameters ($M_{200}$, concentration $c$) per galaxy.

\subsection{Baryonic Tully--Fisher relation}

From Eq.~\eqref{eq:V0}, for universal $k_*$:
\begin{equation}
M_{\mathrm{bar}} = \frac{\pi}{2G_N k_*}\,V_0^2.
\label{eq:BTFR_quadratic}
\end{equation}
Since $k_*$ is approximately constant across galaxies, this
is a linear relation $M_{\mathrm{bar}} \propto V_0^2$.
The observed baryonic Tully--Fisher relation
$M_{\mathrm{bar}} \propto V_0^4$~\cite{McGaughBTFR2000}
is recovered in the deep MOND limit where the Verlinde
contribution dominates (Sec.~\ref{sec:a0}).
The running-$G$ contribution gives the correct
intermediate-radius behavior, while the Verlinde contribution
governs the asymptotic regime where
$v^4 = a_0 G_N M_{\mathrm{bar}}$.

\subsection{Radial acceleration relation}

The RAR~\cite{McGaughRAR2016} relates observed acceleration
$g_{\mathrm{obs}}$ to baryonic acceleration $g_{\mathrm{bar}}$
via
\begin{equation}
g_{\mathrm{obs}} = \frac{g_{\mathrm{bar}}}
{1 - e^{-\sqrt{g_{\mathrm{bar}}/g_\dagger}}},
\label{eq:RAR}
\end{equation}
with $g_\dagger = a_0 = (1.20 \pm 0.02) \times 10^{-10}\;
\mathrm{m/s}^2$.  In the $\tauasym$ framework, the
running-$G$ correction at the galactic scale produces
\begin{equation}
g_{\mathrm{obs}} = g_{\mathrm{bar}} +
\frac{2G_N M_{\mathrm{bar}}\, k_*}{\pi r^2},
\label{eq:g_obs}
\end{equation}
which reproduces the deep-MOND behavior
$g_{\mathrm{obs}} \sim \sqrt{a_0\, g_{\mathrm{bar}}}$
when $k_*$ is related to $a_0$ through the cosmological
connection (Sec.~\ref{sec:a0}).

\subsection{Dark matter density profile}

The apparent dark matter density from
Eq.~\eqref{eq:sigma_grav} is
\begin{equation}
\rho_{\mathrm{DM}}(r) = \frac{1}{4\pi G_N}\,
\frac{k_* V_0^2}{r^2} \propto \frac{1}{r^2},
\label{eq:rho_dm}
\end{equation}
which gives an enclosed mass $M_{\mathrm{DM}}(r) \propto r$.
This is a $1/r^2$ density profile, steeper than isothermal
($1/r^2$, coincidentally the same) but distinct from the
NFW profile ($1/r$ at small $r$, $1/r^3$ at large $r$).
The key prediction is the \emph{absence} of a break:
the $1/r^2$ profile extends to arbitrarily large $r$ without
the $1/r^3$ falloff characteristic of NFW halos.  This is
testable with weak lensing data extending to the
virial radius and beyond.  The KiDS-1000 weak lensing
analysis~\cite{Brouwer2021} already extends the RAR by
two decades in acceleration and is consistent with the
absence of a break.

\subsection{Dwarf spheroidal galaxies}

Ghari and Haghi~\cite{GhariHaghi2026} compared Verlinde's
emergent gravity to MOND for 23 Milky Way dwarf spheroidal
satellites and found Verlinde favored over MOND at
$5.2\sigma$.  Since the Verlinde contribution is a subset
of the $\tauasym$ framework's $\Sigma_{\mathrm{DM}}$ (the
volume-law entanglement component), this result supports the
information-theoretic interpretation.

\subsection{Solar system constraints}

At solar system scales ($r \ll r_c$), the running-$G$
correction is negligible: $G(r)/G_N - 1 \sim 2k_* r/\pi
\sim 2 \times 10^{-14}$ at $r = 1\;\mathrm{AU}$, well
below current constraints
$|G_{\mathrm{eff}}/G_N - 1| < 10^{-5}$~\cite{Will2014}.

%=======================================================================
\section{The MOND Scale from the Crooks Theorem}
\label{sec:a0}
%=======================================================================

\subsection{The $a_0$--cosmology coincidence}

The empirical MOND acceleration scale
$a_0 \approx 1.2 \times 10^{-10}\;\mathrm{m/s}^2$ is
numerically close to $cH_0$:
\begin{equation}
\frac{cH_0}{a_0} \approx 5.7,
\end{equation}
with $H_0 = 70\;\mathrm{km/s/Mpc}$.  More precisely,
$a_0 \approx cH_0/(2\pi) \approx cH_0/6$
(these differ by $\sim\!5\%$).  The coincidence has been
noted by Milgrom~\cite{Milgrom2020} and is a key input to
Verlinde's derivation~\cite{Verlinde2016}, where
$a_0 = cH_0/6$.

\subsection{KMS--Crooks argument}

We provide a thermodynamic argument for this connection.
The de~Sitter vacuum satisfies the KMS condition with modular
time periodicity
\begin{equation}
\beta_{\mathrm{dS}} = \frac{2\pi}{H_0}.
\label{eq:beta_dS}
\end{equation}
The spatial thermal wavelength is
\begin{equation}
\lambda_{\mathrm{dS}} = c\,\beta_{\mathrm{dS}} =
\frac{2\pi c}{H_0},
\label{eq:lambda_dS}
\end{equation}
and the associated acceleration scale is
\begin{equation}
a_0 = \frac{c^2}{\lambda_{\mathrm{dS}}} =
\frac{cH_0}{2\pi} \approx 1.08 \times 10^{-10}\;
\mathrm{m/s}^2.
\label{eq:a0}
\end{equation}

The Crooks fluctuation theorem~\cite{Crooks1999,Kwon2019}
gives the ratio of forward to reverse transition probabilities
as $P_F(\Sigma)/P_R(-\Sigma) = e^\Sigma$.  For a gravitational
process, the entropy production per degree of freedom is
\begin{equation}
\Sigma = \frac{\kB T_U}{\kB T_{\mathrm{dS}}} =
\frac{a}{cH_0},
\label{eq:sigma_a}
\end{equation}
where $T_U = \hbar a/(2\pi c \kB)$ is the Unruh temperature
at acceleration $a$ and
$T_{\mathrm{dS}} = \hbar H_0/(2\pi \kB)$.

Three regimes emerge:
\begin{itemize}
\item $a \gg a_0$ ($\Sigma \gg 1$):
Gravitational processes are irreversible.
The forward process dominates by a factor $e^\Sigma \gg 1$.
The time arrow is strong.  This is the Newtonian regime.

\item $a \sim a_0$ ($\Sigma \sim 1$):
Forward and reverse processes have comparable probabilities
($\sim\!63\%$ vs.\ $\sim\!37\%$).  Irreversibility is
\emph{marginal}.  This is the MOND transition.

\item $a \ll a_0$ ($\Sigma \ll 1$):
Processes are nearly reversible.
The time arrow weakens.  This is the deep MOND regime,
where the de~Sitter thermal bath dominates.
\end{itemize}

The MOND transition is therefore the transition from
strong to marginal irreversibility in the gravitational
channel, as measured by the Crooks entropy
production~\cite{BassoCeleri2025}.

\subsection{Crooks interpolation function}

The exponential form of the Crooks ratio suggests a natural
MOND interpolation function.  Writing $x = g_{\mathrm{bar}}/a_0$,
we propose
\begin{equation}
\mu(x) = 1 - e^{-x}.
\label{eq:mu_crooks}
\end{equation}
The effective acceleration is
$g_{\mathrm{obs}} = g_{\mathrm{bar}}/\mu(x)$.
Properties of Eq.~\eqref{eq:mu_crooks}:
\begin{align}
x \gg 1:&\quad \mu \to 1 \quad (\text{Newtonian}),
\nonumber\\
x \ll 1:&\quad \mu \approx x \quad (\text{deep MOND}).
\nonumber
\end{align}
This is the same asymptotic behavior as the standard MOND
interpolation functions.  Numerically, for representative
values:
\begin{align}
x = 0.1:&\quad \mu_{\mathrm{Crooks}} = 0.095,\quad
\mu_{\mathrm{std}} = 0.091,
\nonumber\\
x = 1.0:&\quad \mu_{\mathrm{Crooks}} = 0.632,\quad
\mu_{\mathrm{std}} = 0.500,
\nonumber\\
x = 3.0:&\quad \mu_{\mathrm{Crooks}} = 0.950,\quad
\mu_{\mathrm{std}} = 0.750.
\nonumber
\end{align}
The Crooks form transitions faster from MOND to Newton than
the standard $\mu(x) = x/(1+x)$; whether this is preferred by
the data requires systematic fitting to SPARC galaxies, which
we leave to future work.

\subsection{Limitations of the $a_0$ argument}

The identification of $a_0 = cH_0/(2\pi)$ should be classified as
a \emph{dimensional argument with physical content}, not a
rigorous derivation.  The scale $cH_0$ follows inevitably from
combining the speed of light with the Hubble rate.  The factor
$2\pi$ is motivated by the KMS periodicity of the de~Sitter
thermal state but is not uniquely derived from the Crooks
theorem.  Verlinde's factor of 6 arises from a different
geometric argument (the volume-to-area ratio of a sphere
combined with the quadratic elastic energy).  Since
$2\pi \approx 6.28 \approx 6$ to within 5\%, these
predictions are observationally indistinguishable with current
data.

The situation is analogous to the Chandrasekhar mass:
$M_{\mathrm{Ch}} \sim (\hbar c/G)^{3/2}/m_p^2$ follows from
dimensional analysis, but the coefficient 0.77 requires
solving the Lane--Emden equation.  Similarly, the precise
$O(1)$ prefactor in $a_0$ awaits a more detailed calculation
of the Crooks entropy production for test particles in
de~Sitter spacetime, using the formalism of
Ref.~\cite{BassoCeleri2025}.

%=======================================================================
\section{The Bullet Cluster}
\label{sec:bullet}
%=======================================================================

The Bullet Cluster (1E\,0657-558)~\cite{Clowe2006} poses the
sharpest challenge to any theory that explains dark matter
through modified gravity or emergent effects.  The lensing
mass peaks are offset from the X-ray gas (which contains
$\sim\!85\%$ of the baryonic mass) by $\sim\!150\;\mathrm{kpc}$,
coinciding instead with the galaxy distribution.

\subsection{The challenge for modified gravity}

For a theory where the gravitational modification is
proportional to the local baryonic mass (as in running $G$),
the gas dominates the total gravitational signal.  Running $G$
cannot by itself produce a lensing peak offset from the gas.
This is a generic problem shared by all modified gravity
theories that have been tested against the Bullet
Cluster~\cite{Famaey2012}, with the exception of
MOG/STVG~\cite{Brownstein2007}, whose additional vector field
provides the needed offset.

\subsection{Entanglement rigidity argument}

We offer a qualitative argument based on the dynamical
timescales.  The volume-law entanglement pattern that
produces the apparent dark matter (the Verlinde component of
$\Sigma_{\mathrm{DM}}$) was established over the equilibration
timescale
$t_{\mathrm{eq}} \sim H_0^{-1} \sim 10^{10}\;\mathrm{yr}$.
The cluster collision timescale is
$t_{\mathrm{coll}} \sim 10^{8}\;\mathrm{yr}$.  The ratio
\begin{equation}
\frac{t_{\mathrm{eq}}}{t_{\mathrm{coll}}} \sim 100
\label{eq:timescale}
\end{equation}
suggests that the entanglement pattern responds sluggishly
to the rapid collision.  During the merger, the gas is
ram-pressure stripped and decelerates, but the entanglement
displacement pattern---associated with the pre-collision mass
distribution---cannot fully re-equilibrate.  If $\sim\!95\%$ of
the entanglement-based $\Sigma_{\mathrm{DM}}$ is ``frozen''
during the collision, the lensing signal would partially track
the pre-collision mass distribution, which is centered on the
galaxy positions.

This argument is qualitative and speculative.  It requires a
theory of entanglement equilibration timescales in de~Sitter
space that does not currently exist.  We note that the Bullet
Cluster is also problematic for $\Lambda$CDM: the inferred
collision velocity ($\sim\!3000\;\mathrm{km/s}$) is
statistically improbable~\cite{Lee2010}, and the Abell~520
cluster shows the \emph{opposite} pattern (a lensing peak
without galaxies), which CDM also struggles to
explain~\cite{Jee2007}.

%=======================================================================
\section{Cosmological Scale: Dark Energy}
\label{sec:cosmo}
%=======================================================================

At cosmological scales, the entropy production from the
de~Sitter vacuum gives
\begin{equation}
\Sigma_{\mathrm{dS}}(r) = -\ln\!\left(1 -
\frac{H_0^2 r^2}{c^2}\right),
\label{eq:sigma_dS}
\end{equation}
which diverges at the Hubble horizon $r_H = c/H_0$,
corresponding to $\tauasym \to 1$: complete information loss
at the cosmological horizon.  This is the ultimate
gravitational $\tauasym$.

Padmanabhan's CosMIn framework~\cite{Padmanabhan2017}
establishes that finite cosmic information requires
$\Lambda > 0$.  In the $\tauasym$ framework, this translates
to: a finite total information loss ($\int \Sigma\,dt < \infty$)
over the history of the universe requires a cosmic information
horizon where $\tauasym = 1$.  This is precisely the
de~Sitter horizon set by $\Lambda$.

The connection between the cosmological constant and the
de~Sitter thermal state provides the bridge to dark matter:
the \emph{same} vacuum entanglement that produces
$\Lambda > 0$ (through the volume-law contribution) also
produces the apparent dark matter at galactic scales (through
its displacement by baryonic matter).  In the
$\Sigma$ language:
\begin{equation}
\Sigma_{\mathrm{total}}(r) \approx
\underbrace{\frac{r_s}{r}}_{\text{Newtonian}} +
\underbrace{\alpha\,\ln\!\left(\frac{r}{r_c}\right)
}_{\text{running }G} +
\underbrace{\beta\,\frac{r}{L_{\mathrm{dS}}}
}_{\text{de Sitter}},
\label{eq:sigma_total}
\end{equation}
where $L_{\mathrm{dS}} = c/H_0$.  Each term dominates in a
different regime: the Newtonian term near the mass source, the
running-$G$ term at galactic scales, and the de~Sitter term at
cosmological scales.  All three emerge as limits of the single
quantum relative entropy
$\Sigma = \KL(\rho_{\mathrm{spacetime}} \|
\rho_{\mathrm{matter}})$~\cite{Paper2,Bianconi2025}.

The DESI DR2 results showing $w(z) \neq -1$ at
$\sim\!4\sigma$~\cite{DESI2024} are naturally compatible with
the $\tauasym$ framework, since a scale-dependent
$\Sigma_{\mathrm{grav}}$ generically produces evolving dark
energy.  We note this as a qualitative consistency, not a
prediction.

%=======================================================================
\section{The CMB Gap}
\label{sec:cmb}
%=======================================================================

The $\tauasym$ framework, as currently formulated,
\emph{cannot} explain the CMB acoustic peaks without
introducing additional degrees of freedom.  This section
presents the argument in full, because the inability to solve
this problem is as informative as the successes.

\subsection{The quantitative failure}

The CMB third acoustic peak requires a pressureless
($c_s^2 = 0$), photon-decoupled component with
$\Omega h^2 \approx 0.12$.  Without such a component,
the predicted first peak position shifts from $\ell \approx 220$
to $\ell \approx 170$ (a 23\% error, $>100\sigma$),
the third-to-second peak ratio drops from 0.97 to $\sim\!0.55$
(a 43\% error, $>20\sigma$), and the matter-radiation equality
redshift shifts from $z_{\mathrm{eq}} \approx 3400$ to
$z_{\mathrm{eq}} \approx 530$ (a factor 6.4
discrepancy)~\cite{Planck2018,Hu2008}.  These are catastrophic
failures, not marginal discrepancies.

\subsection{Why running $G$ is irrelevant at CMB scales}

The IR running formula $G(k) = G_N[1 + k_*/k]$ with
$k_* \sim 10^{-2}\;\mathrm{kpc}^{-1}$ gives $G/G_N \sim 2700$
at CMB scales ($k_{\mathrm{CMB}} \sim 10^{-2}\;
\mathrm{Mpc}^{-1}$), which is physically absurd.  The running
is a late-time, galactic-scale effect that does not apply at
$z \sim 1100$.  Constraints from Planck require
$|G(z\!\sim\!1100)/G_0 - 1| < 0.05$ at 95\% CL.

\subsection{The fundamental obstruction}

The core problem is that any perturbation mode derived from a
local, Lorentz-invariant quantum field theory on an FRW
background has a sound speed $c_s > 0$.  The quantum relative
entropy $\KL(\rho_{\mathrm{spacetime}} \|
\rho_{\mathrm{matter}})$ is Lorentz invariant, so its
perturbations propagate at the speed of light.  But the CMB
requires $c_s = 0$---a non-propagating mode.

The only known ways to achieve $c_s = 0$ are: (i)~breaking
Lorentz invariance in the gravity sector (as in AeST's
aether~\cite{Skordis2021} and the Khronon
field~\cite{Blanchet2024}), (ii)~introducing massive
non-relativistic particles (CDM), or (iii)~imposing a
non-propagating constraint via a Lagrange multiplier (as in
Bianconi's $G$-field~\cite{Bianconi2025}).

\subsection{Structural parallel with AeST and the Modular
Khronon}

The parallel between the $\tauasym$ framework and the
Khronon/AeST theories is notable
(Table~\ref{tab:parallel}):

\begin{table}[t]
\caption{\label{tab:parallel}Structural parallel between the
$\tauasym$ framework and relativistic MOND theories.  The
observer's 4-velocity $u^\mu$ plays the same structural role
as the aether field $A^\mu$.}
\begin{ruledtabular}
\begin{tabular}{lcc}
& \textbf{$\tauasym$ framework} & \textbf{AeST/Khronon} \\
\colrule
Scalar field & $\Sigma$ (entropy prod.) &
$\tau$ (Khronon) \\
Background & $\Sigma = 0$ (comoving) &
$Q = 1$ (comoving) \\
Observer/aether & $u^\mu$ (observer) &
$A^\mu$ (aether) \\
MOND limit & Running $G$ &
$J(Y)$ function \\
CMB behavior & Running $\mu$, dust-like at $z > 100$ &
GDM: $w\!\sim\!0$, $c_s\!\sim\!0$ \\
\end{tabular}
\end{ruledtabular}
\end{table}

In the Khronon theory~\cite{Blanchet2024}, the DBI-inspired
kinetic function $K(Q) = \mu^2(Q-1)^2$ has a minimum at $Q = 1$
(unit timelike gradient), producing a mode with $\omega = 0$
(non-propagating) on the Minkowski background.  This is the
mechanism that yields $c_s = 0$ and hence CDM-like behavior.

In the $\tauasym$ framework, the observer's 4-velocity $u^\mu$
plays the role of the Khronon gradient
$\nabla^\mu \tauasym$.  If $u^\mu$ is promoted from an
external input to a dynamical field---constrained by
$g_{\mu\nu} u^\mu u^\nu = -1$ (unit timelike)---the resulting
theory would naturally contain:
(i)~a preferred foliation (breaking Lorentz invariance in the
gravity sector),
(ii)~a scalar mode with $c_s = 0$ from the unit-timelike
constraint, and
(iii)~a connection to the AeST/Khronon field content.

This ``observer-as-aether'' identification
is a natural route to closing the CMB gap.  It changes
the framework fundamentally: the observer is no longer external
but becomes part of the dynamics.  If successful, dark matter at
cosmological scales would be \emph{the energetic cost of
maintaining a time direction in the universe}.  Paper~IV
develops this idea: the modular Khronon
$\delta T_{\mathrm{mod}} = -2\Phi/H$ emerges from QRE
perturbation theory and has $c_s = 0$ in the matter era,
providing exactly the CDM-like perturbation behavior needed
for the CMB~\cite{Paper4}.

\subsection{The CMB gap in context}

The CMB gap is shared by \emph{all} modified gravity
approaches.  MOND~\cite{Milgrom1983} has no CMB prediction
without a relativistic completion.
Verlinde~\cite{Verlinde2016} has no perturbation theory.
The only existing modified gravity theory that fits both
galactic rotation curves and the CMB is
AeST~\cite{Skordis2021}, and it achieves this by introducing
precisely the field content (vector $+$ scalar) that our
analysis identifies as structurally necessary.  The $\tauasym$
framework's contribution is to identify \emph{why} this field
content is needed: to break Lorentz invariance in a way that
produces a $c_s = 0$ mode.

\subsection{Resolution via running Khronon}
\label{sec:khronon_resolution}

The CMB gap identified above has a candidate resolution within
the $\tauasym$ framework once the Khronon coupling $\mu$ is
promoted from a free parameter to a dimensionally-determined
quantity.

\emph{Dimensional uniqueness of $\mu_0$.}---The Khronon kinetic
function $K(Q) = \mu^2(Q - 1)^2$~\cite{Blanchet2024} requires
a coupling $\mu$ with dimensions of inverse length.  In
de~Sitter space, the only available scales are $H_0$, $c$, $\hbar$,
and $\kB$.  The unique combination with dimensions of inverse
length that is independent of $\hbar$ (as befits a classical
coupling) is
\begin{equation}
\mu_0 = \frac{H_0}{c}
= \frac{2\pi \kB T_{\mathrm{dS}}}{\hbar c},
\label{eq:mu0}
\end{equation}
where $T_{\mathrm{dS}} = \hbar H_0/(2\pi \kB)$ is the
de~Sitter temperature.  This is not an assumption but a
consequence of dimensional analysis: the same thermal scale
that determines $a_0 = cH_0/(2\pi)$ also determines
$\mu_0$.

\emph{Running of $\mu$.}---The anomalous dimension of $\mu$
follows from the same DPI~$+$~extensivity argument that
fixes $\eta_G = 1$ for the gravitational coupling
(Sec.~\ref{sec:galactic}):
\begin{equation}
\eta_\mu = 1, \qquad
\mu(k) = \mu_0\!\left(\frac{k_*}{k}\right),
\label{eq:mu_running}
\end{equation}
because the Khronon field couples to gravity through the same
renormalization group flow.  The DPI forbids $\eta_\mu > 1$
(would violate the de~Sitter entropy bound on the Khronon
sector), while $\eta_\mu < 1$ makes the Khronon IR-irrelevant,
contradicting its required role at CMB scales.
$\eta_\mu = 1$ is the unique marginal value.

\emph{Dust-like behavior at CMB scales.}---On an FRW
background, the Khronon gradient satisfies $Q = 1$ (comoving
frame), and the background coupling runs as
\begin{equation}
\mu_{\mathrm{bg}}(a) = \frac{H(a)}{c},
\label{eq:mu_bg}
\end{equation}
where $H(a)$ is the Hubble parameter at scale factor~$a$.
The background Khronon field amplitude
$\delta(a) \equiv Q_0(a) - 1 = I_0/[2\mu_{\mathrm{bg}}(a)^2]$
decreases at early times because $\mu_{\mathrm{bg}}$ grows
with $H(a)$.  The effective GDM equation-of-state
parameter~\cite{Blanchet2024} is
\begin{equation}
\tilde{w}(z) = \frac{\delta(z)}{2}
= \frac{\delta_0}{2}\!\left(\frac{H_0}{H(z)}\right)^{\!2},
\label{eq:w_tilde}
\end{equation}
where $\delta_0 = 0.340$ at $z = 0$
(fixed by $\Omega_{\mathrm{cdm}} = 0.265$).
At recombination ($z \approx 1100$),
$H(z)/H_0 \approx 2.35 \times 10^4$, giving
\begin{equation}
\tilde{w}(z \approx 1100) \approx 3 \times 10^{-10}
\ll 1.
\label{eq:w_dust}
\end{equation}
The effective sound speed, given by the
Jeans-suppressed formula~\cite{Blanchet2024},
$c_s^2 = \tilde{w}/[1 + \tilde{w}\,c^2 k^2/(4\pi G a^2
\rho_K)]$, satisfies $c_s^2 \lesssim 3 \times 10^{-10}$
at all Planck wavenumbers
($k \in [10^{-3},\,0.3]\;\mathrm{Mpc}^{-1}$),
well below the Thomas--Kopp--Skordis~\cite{TKS2016}
constraint $c_s^2 < 3.4 \times 10^{-6}$.
This confirms the Khronon perturbations behave as dust at
CMB scales, precisely the $c_s = 0$ mode identified as
structurally necessary in Sec.~\ref{sec:cmb}.

\emph{Unification.}---The de~Sitter temperature
$T_{\mathrm{dS}}$ now serves as the single origin of both
galactic and cosmological dark-matter phenomenology:
\begin{equation}
T_{\mathrm{dS}} \;\longrightarrow\;
\begin{cases}
a_0 = cH_0/(2\pi) & \text{(galactic: MOND scale)}, \\
\mu_0 = H_0/c & \text{(CMB: Khronon coupling)}.
\end{cases}
\label{eq:unification}
\end{equation}
At galactic scales, the running of $G$ with $\eta_G = 1$
provides MOND-like rotation curves
(Sec.~\ref{sec:galactic}).  At CMB scales, the running of
$\mu$ with $\eta_\mu = 1$ provides a candidate for the
pressureless, non-propagating mode needed for the acoustic
peaks.  Both originate from the same de~Sitter vacuum
entanglement.
Paper~IV~\cite{Paper4} develops the full cosmological
perturbation theory; here we note only that the mechanism
exists and produces the correct qualitative behavior.

\emph{Status: experimentally unverified.}---The running
Khronon resolution has not been tested against the full Planck
CMB power spectrum.  The qualitative argument (dust-like
$\tilde{w}_0 \ll 1$) is encouraging but does not constitute a
quantitative fit.  Until the numerical computation of
Paper~IV reproduces the observed $C_\ell^{TT}$ spectrum to
within Planck error bars, this resolution should be regarded
as a conjecture.

\emph{Limitations.}---The $\Omega_{\mathrm{DM}} h^2$
prefactor---the precise numerical value that Planck measures
as $0.120 \pm 0.001$---is \emph{not} determined by the
dimensional argument alone.  It depends on the full
nonlinear dynamics of the Khronon field on an FRW background,
which requires numerical integration beyond the scope of this
paper.  We predict the dark matter density to order of
magnitude (from $T_{\mathrm{dS}}$ and $\mu_0$); the exact
prefactor requires the numerical computation presented in
Paper~IV.

%=======================================================================
\section{Testable Predictions}
\label{sec:predictions}
%=======================================================================

The $\tauasym$ framework makes five testable predictions
distinguishable from $\Lambda$CDM and standard MOND:

\emph{1. DM density profile $\propto 1/r^2$ without NFW
break.}---The running-$G$ correction produces an enclosed mass
$M_{\mathrm{DM}}(r) \propto r$ at large $r$, without the
$M_{\mathrm{DM}} \to \mathrm{const}$ behavior of NFW halos.
For a Milky Way--mass galaxy ($M_{\mathrm{bar}} \approx
5 \times 10^{10}\;\Msun$), the predicted dark matter density
at $r = 100\;\mathrm{kpc}$ is
$\rho_{\mathrm{DM}} \approx 3 \times 10^{-4}\;\Msun/
\mathrm{pc}^3$, declining as $r^{-2}$ with no steepening.
Weak lensing profiles extending beyond the virial radius should
show no turnover.  The BIG-SPARC
survey~\cite{Haubner2024} with $\sim\!4000$ galaxies can test
this.

\emph{2. Universal crossover scale and running $G$.}---The
crossover momentum
$k_* \approx 2.7 \times 10^{-2}\;\mathrm{kpc}^{-1}$
(crossover radius $r_c \approx 37\;\mathrm{kpc}$) should be
independent of galaxy mass (within observational scatter).
Concretely, $G(r)/G_N = 1 + 2k_* r/\pi$ predicts:
$G/G_N \approx 1.02$ at $r = 1\;\mathrm{kpc}$,
$G/G_N \approx 1.17$ at $r = 10\;\mathrm{kpc}$, and
$G/G_N \approx 2.72$ at $r = 100\;\mathrm{kpc}$.
A systematic variation $k_*(M_{\mathrm{bar}})$ would indicate
environment-dependent running rather than a universal
fundamental scale.

\emph{3. Interpolation function $\mu(x) = 1 - e^{-x}$.}---The
Crooks-inspired interpolation transitions faster from MOND to
Newtonian than the standard $\mu(x) = x/(1+x)$.  At
$x = 1$, $\mu_{\mathrm{Crooks}} = 0.63$ vs.\
$\mu_{\mathrm{std}} = 0.50$.  This is distinguishable with
high-quality rotation curves near the transition radius.

\emph{4. Observer-dependent $a_0$.}---Since
$\Sigma$ is observer-dependent~\cite{BassoCeleri2025}, the
MOND acceleration scale should depend on the observer's
worldline.  Comoving observers see the standard
$a_0 \approx cH_0/(2\pi) \approx 1.08 \times 10^{-10}\;
\mathrm{m/s}^2$, while observers near massive
objects (galaxy clusters, black holes) should see a locally
modified $a_0$ due to gravitational redshift.  Specifically,
at gravitational potential $\Phi$:
\begin{equation}
a_0^{\mathrm{local}} = a_0\,\sqrt{-g_{00}}
\approx a_0\left(1 + \frac{\Phi}{c^2}\right).
\label{eq:a0_local}
\end{equation}
For galaxies in a cluster with $\Phi/c^2 \sim -10^{-5}$,
the predicted shift is
$\delta a_0/a_0 \sim 10^{-5}$---small but in principle
measurable through precision RAR studies of cluster satellite
galaxies compared to field galaxies.

\emph{5. Lensing vs.\ kinematic DM mass may differ.}---If
$\Sigma_{\mathrm{DM}}$ has both running-$G$ and Verlinde
components with different radial profiles, the
``dark matter mass'' inferred from lensing (which probes
$g_{00} + g_{rr}$) may differ from that inferred from
kinematics (which probes $g_{00}$ alone).  In GR, these are
identical for static metrics; in the $\tauasym$ framework,
the logarithmic correction modifies both $g_{00}$ and
$g_{rr}$ equally, preserving the GR prediction.
However, at the transition between the running-$G$ and
Verlinde regimes, deviations are possible.

%=======================================================================
\section{Discussion}
\label{sec:discussion}
%=======================================================================

\subsection{What is genuinely new}

Six elements distinguish this work from previous
approaches to the dark matter problem:

(i)~The identification of ``dark matter'' with
observer-dependent entropy production
$\Sigma_{\mathrm{DM}} = \Sigma_{\mathrm{cl}} -
\Sigma_{\mathrm{q}}$, which is guaranteed non-negative by
the DPI (Theorem~\ref{thm:dm}).  This is not a model
assumption but a mathematical consequence of the framework.

(ii)~The logical chain connecting the Casini--Huerta
entropic $c$-theorem~\cite{CasiniHuerta2012,CasiniHuerta2017},
the Dorau--Much QRE-to-Einstein link~\cite{DorauMuch2025},
and the Kumar running $G$~\cite{Kumar2025} into a coherent
information-theoretic narrative for galactic rotation curves.
Each step is established; the synthesis is new.

(iii)~The identification of $a_0 \approx cH_0/(2\pi)$ with the
KMS thermal periodicity of de~Sitter space through the Crooks
theorem, interpreting the MOND transition as marginal
irreversibility (Sec.~\ref{sec:a0}; note that the $O(1)$
prefactor is a dimensional argument, not a rigorous
derivation).

(iv)~The Crooks interpolation function $\mu(x) = 1 - e^{-x}$,
which is a specific, falsifiable prediction.

(v)~The identification of the observer's 4-velocity with the
aether/Khronon field, providing a concrete bridge between the
$\tauasym$ framework and relativistic MOND theories.

(vi)~The identification of the Khronon coupling
$\mu_0 = H_0/c$ from dimensional uniqueness, and $\eta_\mu = 1$
from DPI~$+$~extensivity, yielding dust-like behavior at CMB
scales without introducing new free parameters
(Sec.~\ref{sec:khronon_resolution}).  This CMB resolution
is not yet verified against the full Planck power spectrum;
the quantitative test requires Paper~IV.

\subsection{What is not claimed}

We do \emph{not} claim: (i)~that the CMB acoustic peaks
are quantitatively fitted---the running Khronon coupling
$\mu_0 = H_0/c$ with $\eta_\mu = 1$ provides CDM-like
behavior at CMB scales
(Sec.~\ref{sec:khronon_resolution}), but the exact
$\Omega_{\mathrm{DM}} h^2$ prefactor remains undetermined
without the full numerical computation of Paper~IV;
(ii)~that the Bullet Cluster is quantitatively fitted (only
a qualitative argument is given); or (iii)~superiority over
$\Lambda$CDM at all scales.
We \emph{do} note that $\eta = 1$ is derived, not assumed:
Kumar~\cite{Kumar2025} identified $\eta = 1$ as the unique
marginal value from QFT first principles.  Within the
$\tauasym$ framework, this is independently reproduced by
the DPI applied to the gravitational RG flow, combined with
the extensivity of de~Sitter entanglement entropy (volume-law
scaling at cosmological scales~\cite{Verlinde2016}), which
uniquely selects $\eta = 1$ in $3+1$ dimensions.
The argument is: $\eta > 1$ violates the de~Sitter entropy
bound; $\eta < 1$ contradicts extensivity (the correction
becomes IR-irrelevant); $\eta = 1$ is the unique marginal
value producing the logarithmic potential.
In general $d$ dimensions, $\eta = d - 3$.
The same DPI~$+$~extensivity argument fixes
$\eta_\mu = 1$ for the Khronon coupling
(Sec.~\ref{sec:khronon_resolution}), providing the
mechanism---though not yet the precise numerical
coefficient---for CMB-scale dark matter behavior.

\subsection{The three-term $\Sigma$ as a unified equation}

Equation~\eqref{eq:sigma_total} approximates the single
quantum relative entropy
$\Sigma = \KL(\rho_{\mathrm{spacetime}} \|
\rho_{\mathrm{matter}})$ in three regimes.  In the vision of
the four-paper series~\cite{Paper1,Paper2,Paper5}, different
boundary conditions on this one QRE give:
Paper~I ($\Sigma = 0$): flat spacetime, perfect recovery;
Paper~II ($\Sigma = r_s/r$): Schwarzschild, the time arrow
of gravity;
Paper~III ($\Sigma = r_s/r + \alpha\ln(r/r_c)$):
galactic scales, dark matter;
Paper~IV ($\Sigma$ with cosmological boundary conditions):
dark energy and the CMB.
These are not four separate theories but four limits of one
equation, analogous to $F = ma$ producing planetary orbits,
spring oscillations, and fluid dynamics depending on
boundary conditions.

\subsection{The CMB gap as a signpost}

The failure at CMB scales is not a deficiency to be
papered over but a signpost.  It points to a specific
mathematical program: determine whether the constrained
minimization of
$\KL(\rho_{\mathrm{spacetime}} \| \rho_{\mathrm{matter}})$
produces Lagrange multiplier fields with $c_s = 0$.  If it
does, the $\tauasym$ framework would derive CDM-like behavior
from quantum information theory.  If it does not, then CDM
particles are a necessary ingredient that the framework must
incorporate rather than explain.  Either outcome would
clarify the scope of the framework.

The running Khronon analysis of Sec.~\ref{sec:khronon_resolution}
suggests an affirmative answer: the constrained observer
dynamics produces a dust-like mode at CMB scales, via the
running coupling $\mu_0 = H_0/c$.  Whether this survives a
full numerical comparison with the Planck power spectrum
remains to be determined (Paper~IV).

\subsection{Division of labor: $G$ and $\mu$}

The $\tauasym$ framework now exhibits a clear division of
labor between two running couplings, both originating from the
de~Sitter vacuum and both with anomalous dimension
$\eta = 1$:

\emph{Running $G$} governs galactic scales.  The infrared
running $G(k) = G_N(1 + k_*/k)$ produces the logarithmic
correction to the Newtonian potential, flat rotation curves,
and the radial acceleration relation.  The crossover scale
$k_* \sim 10^{-2}\;\mathrm{kpc}^{-1}$ is set by
$a_0 = cH_0/(2\pi)$.  This is a modification of the
\emph{strength} of gravity.

\emph{Running $\mu$} governs cosmological scales.  The
Khronon coupling $\mu(k) = \mu_0(k_*/k)$ with
$\mu_0 = H_0/c$ determines the dynamics of the observer
field.  At CMB scales, $\mu_{\mathrm{bg}}(a) = H(a)/c$
produces a pressureless mode ($\tilde{w}_0 \ll 1$) that
mimics cold dark matter perturbations.  This is a
modification of the \emph{observer dynamics}---the energetic
cost of maintaining a time direction.

The two effects are complementary, not redundant: running
$G$ is a late-time, small-scale effect that is irrelevant at
$z \sim 1100$ (Sec.~\ref{sec:cmb}), while running $\mu$ is
a large-scale effect whose galactic contribution is
subdominant to the running-$G$ correction.  Together, they
cover the full range from kiloparsec to gigaparsec scales
with a single origin: the de~Sitter temperature
$T_{\mathrm{dS}}$.

\begin{acknowledgments}
The author thanks M.~M.~Wilde for valuable feedback on
the $\tauasym$ framework, and A.~J.~Parzygnat and
F.~Buscemi for foundational work on retrodiction.
\end{acknowledgments}

%=======================================================================
% BIBLIOGRAPHY
%=======================================================================

\begin{thebibliography}{60}

\bibitem{Rubin1980}
V.~C. Rubin, W.~K. Ford, Jr., and N.~Thonnard,
``Rotational properties of 21 Sc galaxies with a large range
of luminosities and radii, from NGC 4605 ($R = 4\;\mathrm{kpc}$)
to UGC 2885 ($R = 122\;\mathrm{kpc}$),''
Astrophys.\ J.\ \textbf{238}, 471 (1980).

\bibitem{Planck2018}
N.~Aghanim \emph{et~al.} (Planck Collaboration),
``Planck 2018 results.\ VI.\ Cosmological parameters,''
Astron.\ Astrophys.\ \textbf{641}, A6 (2020);
arXiv:1807.06209.

\bibitem{Clowe2006}
D.~Clowe, M.~Brada\v{c}, A.~H. Gonzalez, M.~Markevitch,
S.~W. Randall, C.~Jones, and D.~Zaritsky,
``A direct empirical proof of the existence of dark matter,''
Astrophys.\ J.\ Lett.\ \textbf{648}, L109 (2006);
astro-ph/0608407.

\bibitem{Schumann2019}
M.~Schumann,
``Direct detection of WIMP dark matter: Concepts and status,''
J.\ Phys.\ G \textbf{46}, 103003 (2019);
arXiv:1903.03026.

\bibitem{McGaughRAR2016}
S.~S. McGaugh, F.~Lelli, and J.~M. Schombert,
``Radial acceleration relation in rotationally supported
galaxies,''
Phys.\ Rev.\ Lett.\ \textbf{117}, 201101 (2016);
arXiv:1609.05917.

\bibitem{McGaughBTFR2000}
S.~S. McGaugh, J.~M. Schombert, G.~D. Bothun, and
W.~J.~G. de~Blok,
``The baryonic Tully--Fisher relation,''
Astrophys.\ J.\ Lett.\ \textbf{533}, L99 (2000);
astro-ph/0003001.

\bibitem{Milgrom1983}
M.~Milgrom,
``A modification of the Newtonian dynamics as a possible
alternative to the hidden mass hypothesis,''
Astrophys.\ J.\ \textbf{270}, 365 (1983).

\bibitem{Skordis2021}
C.~Skordis and T.~Zlo\v{s}nik,
``New relativistic theory for modified Newtonian dynamics,''
Phys.\ Rev.\ Lett.\ \textbf{127}, 161302 (2021);
arXiv:2007.00082.

\bibitem{Verlinde2016}
E.~P. Verlinde,
``Emergent gravity and the dark universe,''
SciPost Phys.\ \textbf{2}, 016 (2017);
arXiv:1611.02269.

\bibitem{Ettori2019}
S.~Ettori, V.~Ghirardini, D.~Eckert, E.~Pointecouteau,
F.~Gastaldello, M.~Sereno, M.~Gaspari, S.~Ghizzardi,
M.~Roncarelli, and M.~Rossetti,
``Hydrostatic mass profiles in X-COP galaxy clusters,''
Astron.\ Astrophys.\ \textbf{621}, A39 (2019);
arXiv:1805.00035.

\bibitem{Paper5}
S.-K. Huang,
``Observer-dependent temporal asymmetry: From quantum erasure
to gravitational time'' (Paper~V), in preparation (2026).

\bibitem{Junge2018}
M.~Junge, R.~Renner, D.~Sutter, M.~M. Wilde, and A.~Winter,
``Universal recovery maps and approximate sufficiency of
quantum relative entropy,''
Ann.\ Henri Poincar\'e \textbf{19}, 2955 (2018).

\bibitem{Paper2}
S.-K. Huang,
``Gravitational temporal asymmetry from information loss''
(Paper~II), in preparation (2026).

\bibitem{CasiniHuerta2017}
H.~Casini, E.~Teste, and G.~Torroba,
``Relative entropy and the RG flow,''
J.\ High Energy Phys.\ \textbf{03}, 089 (2017);
arXiv:1611.00016.

\bibitem{CasiniHuerta2012}
H.~Casini and M.~Huerta,
``On the RG running of the entanglement entropy of a circle,''
Phys.\ Rev.\ D \textbf{85}, 125016 (2012);
arXiv:1202.5650.

\bibitem{DorauMuch2025}
V.~Dorau and A.~Much,
``From quantum relative entropy to the semiclassical Einstein
equations,''
Phys.\ Rev.\ Lett.\ \textbf{136}, 091602 (2026);
arXiv:2510.24491.

\bibitem{Kumar2025}
N.~Kumar,
``Marginal IR running of gravity as a natural explanation
for dark matter,''
Phys.\ Lett.\ B \textbf{871}, 140008 (2025);
arXiv:2509.05246.

\bibitem{Reuter2004}
M.~Reuter and H.~Weyer,
``Running Newton constant, improved gravitational actions, and
galaxy rotation curves,''
Phys.\ Rev.\ D \textbf{70}, 124028 (2004);
hep-th/0410117.

\bibitem{Donoghue1994}
J.~F. Donoghue,
``General relativity as an effective field theory: The leading
quantum corrections,''
Phys.\ Rev.\ D \textbf{50}, 3874 (1994);
gr-qc/9405057.

\bibitem{Gubitosi2024}
G.~Gubitosi, O.~F. Piattella, and L.~Casarini,
``Phenomenology of renormalization group improved gravity from
the kinematics of SPARC galaxies,''
Phys.\ Rev.\ D \textbf{110}, 124014 (2024);
arXiv:2403.00531.

\bibitem{SPARC2016}
F.~Lelli, S.~S. McGaugh, and J.~M. Schombert,
``SPARC: Mass models for 175 disk galaxies with Spitzer
photometry and accurate rotation curves,''
Astron.\ J.\ \textbf{152}, 157 (2016);
arXiv:1606.09251.

\bibitem{Brouwer2021}
M.~M. Brouwer \emph{et~al.},
``The radial acceleration relation in GAMA/KiDS:
Bridging the gap between theory and observations,''
Astron.\ Astrophys.\ \textbf{650}, A113 (2021);
arXiv:2106.11677.

\bibitem{GhariHaghi2026}
A.~Ghari and H.~Haghi,
``Verlinde's emergent gravity vs.\ MOND in dwarf spheroidal
galaxies,''
arXiv:2601.01715 (2026).

\bibitem{Will2014}
C.~M. Will,
``The confrontation between general relativity and experiment,''
Living Rev.\ Relativ.\ \textbf{17}, 4 (2014);
arXiv:1403.7377.

\bibitem{Milgrom2020}
M.~Milgrom,
``The $a_0$--cosmology connection in MOND,''
arXiv:2001.09729 (2020).

\bibitem{Crooks1999}
G.~E. Crooks,
``Entropy production fluctuation theorem and the nonequilibrium
work relation for free energy differences,''
Phys.\ Rev.\ E \textbf{60}, 2721 (1999);
cond-mat/9901352.

\bibitem{Kwon2019}
H.~Kwon and M.~S. Kim,
``Fluctuation theorems for a quantum channel,''
Phys.\ Rev.\ X \textbf{9}, 031029 (2019).

\bibitem{BassoCeleri2025}
M.~L.~W. Basso, J.~Maziero, and L.~C. Celeri,
``Quantum detailed fluctuation theorem in curved spacetimes:
The observer dependent nature of entropy production,''
Phys.\ Rev.\ Lett.\ \textbf{134}, 050406 (2025);
arXiv:2405.03902.

\bibitem{Famaey2012}
B.~Famaey and S.~S. McGaugh,
``Modified Newtonian dynamics (MOND): Observational
phenomenology and relativistic extensions,''
Living Rev.\ Relativ.\ \textbf{15}, 10 (2012);
arXiv:1112.3960.

\bibitem{Brownstein2007}
J.~R. Brownstein and J.~W. Moffat,
``The Bullet Cluster 1E0657-558 evidence shows modified gravity
in the absence of dark matter,''
Mon.\ Not.\ R.\ Astron.\ Soc.\ \textbf{382}, 29 (2007);
astro-ph/0702146.

\bibitem{Lee2010}
J.~Lee and E.~Komatsu,
``Bullet Cluster: A challenge to $\Lambda$CDM cosmology,''
Astrophys.\ J.\ \textbf{718}, 60 (2010);
arXiv:1003.0939.

\bibitem{Jee2007}
M.~J. Jee \emph{et~al.},
``Discovery of a ringlike dark matter structure in the core of
the galaxy cluster Cl\,0024+17,''
Astrophys.\ J.\ \textbf{661}, 728 (2007);
arXiv:0706.3048.

\bibitem{Blanchet2024}
L.~Blanchet and C.~Skordis,
``A new class of relativistic MOND theories,''
J.\ Cosmol.\ Astropart.\ Phys.\ \textbf{11}, 040 (2024);
arXiv:2404.06584.

\bibitem{Bianconi2025}
G.~Bianconi,
``Gravity from entropy,''
Phys.\ Rev.\ D \textbf{111}, 066001 (2025);
arXiv:2408.14391.

\bibitem{Padmanabhan2017}
T.~Padmanabhan,
``Cosmic information, the cosmological constant and the
amplitude of primordial perturbations,''
Phys.\ Lett.\ B \textbf{773}, 81 (2017);
arXiv:1703.06144.

\bibitem{DESI2024}
DESI Collaboration,
``DESI 2024 VI: Cosmological constraints from the measurements
of baryon acoustic oscillations,''
arXiv:2404.03002 (2024).

\bibitem{Hu2008}
W.~Hu,
``Lecture notes on CMB theory: From nucleosynthesis to
recombination,''
arXiv:0802.3688 (2008).

\bibitem{Haubner2024}
M.~Haubner, F.~Lelli, \emph{et~al.},
``BIG-SPARC: A new galaxy sample for testing dark matter and
galaxy formation models,''
arXiv:2411.13329 (2024).

\bibitem{Paper1}
S.-K. Huang,
``The arrow of time from Petz recovery'' (Paper~I),
arXiv:2603.xxxxx (2026).

\bibitem{Jacobson1995}
T.~Jacobson,
``Thermodynamics of spacetime: The Einstein equation of state,''
Phys.\ Rev.\ Lett.\ \textbf{75}, 1260 (1995);
gr-qc/9504004.

\bibitem{Jacobson2016}
T.~Jacobson,
``Entanglement equilibrium and the Einstein equation,''
Phys.\ Rev.\ Lett.\ \textbf{116}, 201101 (2016);
arXiv:1505.04753.

\bibitem{Smolin2017}
L.~Smolin,
``MOND as a regime of quantum gravity,''
Phys.\ Rev.\ D \textbf{96}, 083523 (2017);
arXiv:1704.00780.

\bibitem{Pazy2013}
E.~Pazy,
``Quantum statistical modified entropic gravity as a theoretical
basis for MOND,''
Phys.\ Rev.\ D \textbf{87}, 084063 (2013);
arXiv:1302.4411.

\bibitem{deNetto2022}
T.~de~Paula~Netto, L.~Modesto, and I.~L. Shapiro,
``Universal leading quantum correction to the Newton potential,''
Eur.\ Phys.\ J.\ C \textbf{82}, 160 (2022);
arXiv:2110.14263.

\bibitem{Donoghue2020}
J.~F. Donoghue,
``A critique of the asymptotic safety program,''
Front.\ Phys.\ \textbf{8}, 56 (2020);
arXiv:1911.02967.

\bibitem{Deddo2024}
E.~Deddo, J.~Kang, L.~Ren, and A.~C. Wall,
``Entropic irreversibility theorems,''
J.\ High Energy Phys.\ \textbf{09}, 179 (2024);
arXiv:2404.15077.

\bibitem{Grumiller2010}
D.~Grumiller,
``Model for gravity at large distances,''
Phys.\ Rev.\ Lett.\ \textbf{105}, 211303 (2010);
arXiv:1011.3625.

\bibitem{Brouwer2017}
M.~M. Brouwer \emph{et~al.},
``First test of Verlinde's theory of emergent gravity using
weak gravitational lensing measurements,''
Mon.\ Not.\ R.\ Astron.\ Soc.\ \textbf{466}, 2547 (2017);
arXiv:1612.03034.

\bibitem{DeVuyst2025}
S.~De~Vuyst, P.~A. H\"ohn, A.~Kirklin, and J.~Eccles,
``Gravitational entropy is observer-dependent,''
J.\ High Energy Phys.\ \textbf{07}, 146 (2025).

\bibitem{Parzygnat2023}
A.~J. Parzygnat and F.~Buscemi,
``Axioms for retrodiction: Achieving time-reversal symmetry
with a prior,''
Quantum \textbf{7}, 1013 (2023).

\bibitem{Fawzi2015}
O.~Fawzi and R.~Renner,
``Quantum conditional mutual information and approximate
Markov chains,''
Commun.\ Math.\ Phys.\ \textbf{340}, 575 (2015).

\bibitem{TKS2016}
S.~Thomas, M.~Kopp, and C.~Skordis,
``Constraining generalised dark matter parameters with
the CMB,''
Astrophys.\ J.\ \textbf{830}, 155 (2016);
arXiv:1601.05097.

\bibitem{Paper4}
S.-K.~Huang,
``One equation, three scales:
$\Sigma = D(\rho_{\mathrm{spacetime}} \| \rho_{\mathrm{matter}})$
as the universal action,''
Paper~IV in this series (2026).

\end{thebibliography}

\end{document}
